\section{Metodología}
    El desarrollo se implementó en dos scripts independientes: \textit{Python}
    genera el conjunto de datos, ejecuta las cinco pruebas y produce las figuras;
    \textit{R} lee el mismo archivo y recalcula todo desde cero, de modo que
    cualquier discrepancia delataría un error de implementación. El núcleo
    estadístico de ambos se reproduce en los Anexos \ref{anexo_python} y
    \ref{anexo_r}. Todas las pruebas usan $\alpha = 0{,}05$.

    \subsection{Conjunto de datos}
        Se emplea un conjunto simulado de consumo energético mensual de 300
        clientes, 100 por cada sector, generado con semilla fija (42) para
        garantizar la reproducibilidad. Las variables se resumen en la Tabla
        \ref{dataset_variables}.

        \begin{table}[H]
            \centering
            \caption{Variables del conjunto de datos}
            \begin{center}
                \begin{tabular}{|l|l|>{\raggedright\arraybackslash}p{2.9cm}|}
                    \hline
                    \textbf{Variable} & \textbf{Tipo} & \textbf{Descripción} \\
                    \hline
                    id\_cliente & Nominal & Identificador único
                    (CL-0001 a CL-0300) \\
                    \hline
                    sector & Nominal & Residencial, Comercial o Industrial \\
                    \hline
                    region & Nominal & Andina, Caribe o Pacífica \\
                    \hline
                    temperatura\_c & Continua & Temperatura promedio del mes
                    en \textdegree C \\
                    \hline
                    consumo\_kwh & Continua & Consumo mensual en
                    kilovatios-hora \\
                    \hline
                \end{tabular}
                \label{dataset_variables}
            \end{center}
        \end{table}

        El consumo se construye según la ecuación (\ref{generador}), donde
        $\mu_s$ y $\sigma_s$ son la media y la desviación propias de cada sector y
        $T_i$ la temperatura del cliente, cuya región se asigna con independencia
        del sector.

        \begin{equation}
            Y_i = \mathcal{N}(\mu_s, \sigma_s^2) + 4\,(T_i - 20)
            \label{generador}
        \end{equation}

        El coeficiente 4 es el dato decisivo del experimento: la relación
        temperatura consumo existe por diseño, es idéntica en los tres sectores y
        su magnitud se conoce de antemano, lo que convierte al conjunto en un banco
        de pruebas donde puede verificarse qué recupera cada estadístico. Los
        sectores, en cambio, se separan por niveles muy distintos, con medias de
        180, 420 y 900 kWh y desviaciones de 40, 90 y 180 kWh respectivamente.

    \subsection{Diseño del contraste}
        Las cinco hipótesis se resumen en la Tabla \ref{hipotesis}. Las dos primeras
        no son un trámite previo sino que condicionan a las siguientes, ya que
        determinan si las pruebas paramétricas son aplicables y bajo qué variante.

        \begin{table}[H]
            \centering
            \caption{Hipótesis contrastadas e implementación}
            \begin{center}
                \begin{tabular}{|c|>{\raggedright\arraybackslash}p{2.4cm}|l|l|}
                    \hline
                    \textbf{\#} & \textbf{Hipótesis nula} & \textbf{Python} &
                    \textbf{R} \\
                    \hline
                    1 & El consumo de cada sector es normal & \texttt{shapiro} &
                    \texttt{shapiro.test} \\
                    \hline
                    2 & Las varianzas son homogéneas & \texttt{levene} &
                    \texttt{bartlett.test} \\
                    \hline
                    3 & $\mu_{Res} = \mu_{Com}$ & \texttt{ttest\_ind} &
                    \texttt{t.test} \\
                    \hline
                    4 & $\mu_{Res} = \mu_{Com} = \mu_{Ind}$ & \texttt{anova\_lm} &
                    \texttt{aov} \\
                    \hline
                    5 & $\rho_{T,Y} = 0$ & \texttt{pearsonr} &
                    \texttt{cor.test} \\
                    \hline
                \end{tabular}
                \label{hipotesis}
            \end{center}
        \end{table}

        La comparación de dos medias emplea el estadístico de la ecuación
        (\ref{welch}), que no asume varianzas iguales y ajusta los grados de
        libertad mediante la aproximación de Welch-Satterthwaite
        \cite{welch1947generalization}.

        \begin{equation}
            t = \frac{\bar{x}_1 - \bar{x}_2}
                     {\sqrt{\dfrac{s_1^2}{n_1} + \dfrac{s_2^2}{n_2}}}
            \label{welch}
        \end{equation}

        Para tres o más grupos se recurre al ANOVA de un factor, cuyo estadístico
        de la ecuación (\ref{anova}) contrasta la variabilidad entre grupos frente
        a la que existe dentro de ellos \cite{montgomery2020design}.

        \begin{equation}
            F = \frac{SC_{entre} / (k-1)}{SC_{dentro} / (N-k)}
            \label{anova}
        \end{equation}

        Un $F$ significativo indica que al menos una media difiere, pero no cuál,
        por lo que se aplica el post-hoc de Tukey de la ecuación (\ref{tukey}), que
        controla la tasa de error por familia al examinar las tres comparaciones
        por pares \cite{tukey1949comparing}.

        \begin{equation}
            HSD = q_{\alpha,k,N-k} \sqrt{\frac{CM_{dentro}}{n}}
            \label{tukey}
        \end{equation}

        Por último, la relación entre temperatura y consumo se evalúa con el
        coeficiente de Pearson de la ecuación (\ref{pearson})
        \cite{pearson1895regression} y con la recta de regresión estimada por
        mínimos cuadrados de la ecuación (\ref{ols}).

        \begin{equation}
            r = \frac{\sum (T_i - \bar{T})(Y_i - \bar{Y})}
                     {\sqrt{\sum (T_i - \bar{T})^2 \sum (Y_i - \bar{Y})^2}}
            \label{pearson}
        \end{equation}

        \begin{equation}
            \hat{Y} = b_0 + b_1 T
            \label{ols}
        \end{equation}

        Ambas cantidades se calculan a nivel global y, de forma separada, dentro de
        cada sector. Esa doble mirada es la que permite detectar el efecto de la
        agregación que se analiza en la Sección \ref{seccion_visualizacion}.

    \subsection{Magnitud del efecto y potencia}
        \label{subsec_efecto}
        Un valor $p$ responde si una diferencia es distinguible del azar, no si
        importa; con $n = 100$ por grupo casi cualquier separación alcanza la
        significancia, de modo que reportar únicamente el valor $p$ reproduce el
        problema que advierte la declaración de la ASA \cite{wasserstein2016asa}.
        Cada contraste se acompaña por ello de su tamaño de efecto
        \cite{cohen1988statistical}, \cite{lakens2013calculating}. Para la
        comparación de dos medias se emplea la $d$ de Cohen de la ecuación
        (\ref{cohend}), con $s_p$ la desviación combinada, y su corrección de
        sesgo $g$ de Hedges.

        \begin{equation}
            d = \frac{\bar{x}_1 - \bar{x}_2}{s_p}, \qquad
            s_p = \sqrt{\frac{(n_1-1)s_1^2 + (n_2-1)s_2^2}{n_1 + n_2 - 2}}
            \label{cohend}
        \end{equation}

        Para el ANOVA se reportan la proporción de varianza explicada $\eta^2$ y
        su versión corregida $\omega^2$ de la ecuación (\ref{omega}), que
        descuenta el sesgo optimista de la primera.

        \begin{equation}
            \eta^2 = \frac{SC_{entre}}{SC_{total}}, \qquad
            \omega^2 = \frac{SC_{entre} - (k-1)\,CM_{dentro}}
                            {SC_{total} + CM_{dentro}}
            \label{omega}
        \end{equation}

        La potencia observada se calcula sobre la distribución no central
        correspondiente a cada estadístico. A ella se añade un análisis de
        sensibilidad sobre la quinta hipótesis: la correlación más pequeña
        detectable con potencia $1-\beta = 0{,}80$ y $n$ observaciones se obtiene
        invirtiendo la transformación de Fisher según la ecuación
        (\ref{rdetectable}). Ese umbral separa dos explicaciones posibles de una
        correlación no significativa, la falta de tamaño muestral y la pequeñez
        real del coeficiente.

        \begin{equation}
            r_{\min} = \tanh\!\left(\frac{z_{1-\alpha/2} + z_{1-\beta}}
                                         {\sqrt{n-3}}\right)
            \label{rdetectable}
        \end{equation}

    \subsection{Contrastes robustos a la heterocedasticidad}
        \label{subsec_robustez}
        Tanto el ANOVA clásico como el post-hoc de Tukey asumen varianzas
        iguales, supuesto que la segunda hipótesis rechaza. En lugar de dejar
        constancia de ello como una limitación, todo el bloque de comparación de
        medias se repite con los procedimientos que no lo exigen: el ANOVA de
        Welch de la ecuación (\ref{welchanova}), que pondera cada grupo por
        $w_i = n_i/s_i^2$ \cite{welch1951comparison}, y el post-hoc de
        Games-Howell de la ecuación (\ref{gameshowell}), que sustituye el error
        estándar común de Tukey por el de Welch en cada par y ajusta los grados
        de libertad mediante Welch-Satterthwaite \cite{games1976pairwise}.

        \begin{equation}
            F_W = \frac{\dfrac{1}{k-1}\sum w_i (\bar{x}_i - \bar{x}_w)^2}
                       {1 + \dfrac{2(k-2)}{k^2-1}\,\Lambda}
            \label{welchanova}
        \end{equation}

        \begin{equation}
            q_{ij} = \frac{|\bar{x}_i - \bar{x}_j|}
                          {\sqrt{\dfrac{1}{2}\left(\dfrac{s_i^2}{n_i}
                                 + \dfrac{s_j^2}{n_j}\right)}}
            \label{gameshowell}
        \end{equation}

        Si ambos procedimientos conducen a la misma decisión, la conclusión queda
        respaldada con independencia del supuesto; si divergieran, habría que
        reportar la versión robusta. El resultado de esa comprobación se presenta
        en la Sección \ref{subsec_robustez_resultados}.

    \subsection{Plan de visualización}
        \label{subsec_plan_visual}
        Cada hipótesis contrastada tiene asignada al menos una figura, y cada
        figura lleva impreso el estadístico y el valor $p$ que representa, de modo
        que pueda leerse sin emparejarla con la tabla correspondiente. Las nueve
        figuras se construyen dos veces, en Python y en R, sobre estadísticos
        calculados de forma independiente en cada entorno. La Tabla
        \ref{plan_figuras} resume esa correspondencia entre hipótesis, biblioteca
        y figura.

        \begin{table}[H]
            \centering
            \caption{Correspondencia entre hipótesis, biblioteca y figura}
            \begin{center}
                \begin{tabular}{|c|>{\raggedright\arraybackslash}p{2.05cm}|l|c|}
                    \hline
                    \textbf{\#} & \textbf{Figura} & \textbf{Biblioteca} &
                    \textbf{Ref.} \\
                    \hline
                    1 & Histograma con normal teórica & Matplotlib / ggplot2 &
                    \ref{fig_normalidad} \\
                    \hline
                    1 & Q-Q por sector & Matplotlib / ggplot2 &
                    \ref{fig_qq} \\
                    \hline
                    2 & Violín con caja y puntos & seaborn / ggplot2 &
                    \ref{fig_boxplot} \\
                    \hline
                    3 & Medias con IC del 95\,\% & Matplotlib / ggplot2 &
                    \ref{fig_medias} \\
                    \hline
                    4 & Forest de Tukey y Games-Howell & Matplotlib / ggplot2 &
                    \ref{fig_forest} \\
                    \hline
                    5 & Dispersión global con recta OLS & seaborn / ggplot2 &
                    \ref{fig_regresion} \\
                    \hline
                    5 & Dispersión por sector \textbf{interactiva} &
                    Plotly / ggplotly & \ref{fig_plotly} \\
                    \hline
                    5 & Descomposición de la atenuación & seaborn / ggplot2 &
                    \ref{fig_atenuacion} \\
                    \hline
                    1--5 & \textbf{Dashboard} de cuatro paneles
                    \textbf{interactivo} & Plotly / ggplotly &
                    \ref{fig_dashboard} \\
                    \hline
                \end{tabular}
                \label{plan_figuras}
            \end{center}
        \end{table}

        Las dos figuras interactivas se exportan como HTML autónomo, junto a sus 
        equivalentes de R generados con \texttt{ggplotly}.
        En ellas la interacción cumple una función analítica y no decorativa: aislar
        un sector desde la leyenda equivale a ejecutar el análisis condicionado que
        sostiene el argumento de este informe, y el \textit{hover} identifica
        cliente y región sin volver a los datos.
